\chapter{Especificação Técnica Normativa do Formato Twee 3}
\label{an:especificacao-twee3}

Este anexo apresenta um extrato normativo e resumo técnico da especificação aberta do formato **Twee 3**, publicada e mantida pela \textit{Interactive Fiction Technology Foundation} (\gls{iftf}). A adoção deste formato como base de persistência do \gls{piap} assegura a interoperabilidade com compiladores padrão do ecossistema Twine (como o \textit{Tweego}).

\section{Sintaxe Geral e Passagens}
Um ficheiro Twee 3 é codificado em formato de texto plano UTF-8. A unidade estrutural básica é a \textbf{passagem}. Cada passagem inicia-se com um cabeçalho obrigatório, seguido de metadados opcionais, tags e, nas linhas seguintes, o corpo de texto da passagem.

\subsection{Sintaxe do Cabeçalho da Passagem}
O cabeçalho de uma passagem deve ser escrito numa única linha, obedecendo à seguinte gramática formal:
\begin{verbatim}
:: NomeDaPassagem [tags] {metadados}
\end{verbatim}
\begin{itemize}
    \item \texttt{::} -- Delimitador inicial obrigatório. Deve ser posicionado no início de uma linha sem espaços precedentes.
    \item \texttt{NomeDaPassagem} -- O título identificador único da passagem. Não pode conter os caracteres \texttt{[} ou \texttt{\{} e é sensível a maiúsculas/minúsculas (\textit{case-sensitive}).
    \item \texttt{[tags]} -- Lista opcional de palavras-chave separadas por espaço, delimitada por parênteses retos.
    \item \texttt{\{metadados\}} -- Objeto opcional em formato \gls{json} delimitado por chavetas, contendo propriedades de layout ou de sistema.
\end{itemize}

\section{Especificação de Metadados de Posicionamento}
A especificação Twee 3 padroniza as coordenadas espaciais das passagens para permitir que editores gráficos visualizem o grafo de forma coerente.

\subsection{Campos de Geometria \gls{json}}
Os seguintes campos são reservados dentro do objeto de metadados \gls{json} do cabeçalho:
\begin{itemize}
    \item \texttt{"position"} -- Um objeto contendo as propriedades numéricas de posição absoluta no ecrã:
    \begin{itemize}
        \item \texttt{"x"} -- Coordenada horizontal em píxeis.
        \item \texttt{"y"} -- Coordenada vertical em píxeis.
    \end{itemize}
    \item \texttt{"size"} -- Um objeto opcional definindo as dimensões físicas da caixa no canvas:
    \begin{itemize}
        \item \texttt{"width"} -- Largura da caixa em píxeis.
        \item \texttt{"height"} -- Altura da caixa em píxeis.
    \end{itemize}
\end{itemize}

No \gls{piap}, para permitir o agrupamento em Zonas, estendeu-se esta especificação adicionando a propriedade \texttt{"parentId"} que armazena a string do ID da Zona mãe à qual o nó pertence geometricamente.

\section{Especificação de Hiperligações e Transições}
A especificação Twee 3 define que as arestas do grafo são declaradas de forma textual no corpo das passagens. O parser deve identificar hiperligações expressas em parênteses retos duplos (\texttt{[[...]]}). As formas válidas de declaração são:
\begin{enumerate}
    \item \textbf{Ligação Direta:} Onde o texto exibido coincide com o ID do nó de destino:
    \begin{verbatim}
    [[CenaDestino]]
    \end{verbatim}
    \item \textbf{Ligação com Alias (Sintaxe Clássica):} Onde se separa o texto visível do destino lógico por uma barra vertical:
    \begin{verbatim}
    [[Texto da Escolha|CenaDestino]]
    \end{verbatim}
    \item \textbf{Ligação com Seta Direcionada:} Permite desenhar setas explícitas de navegação lógica:
    \begin{verbatim}
    [[Texto da Escolha->CenaDestino]]
    \end{verbatim}
\end{enumerate}

O analisador léxico \texttt{tweeParser.js} do \gls{piap} implementa expressões regulares robustas para extrair estes três formatos, garantindo compatibilidade bidirecional absoluta com a norma da \gls{iftf}.

Cabe salientar que, embora a norma da \gls{iftf} permita tecnicamente a utilização de variáveis lógicas como destinos dinâmicos de hiperligações (ex: \texttt{[[Ir para o local|\$local]]}), o parser do \gls{piap} restringe esta prática, exigindo que todos os destinos do grafo de história sejam identificadores estáticos de passagens. Esta restrição é uma decisão de engenharia essencial para assegurar que a lista de adjacências permaneça estática e determinista, viabilizando a validação em tempo real e a simulação lógica do espaço de estados sem falhas estruturais.